\documentclass[10pt,a4paper]{amsart}
\usepackage{amsmath,amssymb,amsthm,graphicx}
% For commutative diagrams we use Xy-pic
\usepackage[all]{xy}

% We use 2cell for 2-commutative diagrams.
\xyoption{2cell}
% Theorem environments.
%
\theoremstyle{plain}
\newtheorem{theorem}[subsection]{Theorem}
\newtheorem{proposition}[subsection]{Proposition}
\newtheorem{lemma}[subsection]{Lemma}

\theoremstyle{definition}
\newtheorem{definition}[subsection]{Definition}
\newtheorem{example}[subsection]{Example}
\newtheorem{exercise}[subsection]{Exercise}
\newtheorem{situation}[subsection]{Situation}

\theoremstyle{remark}
\newtheorem{remark}[subsection]{Remark}
\newtheorem{remarks}[subsection]{Remarks}

\numberwithin{equation}{subsection}

% Macros
%
\def\lim{\mathop{\mathrm{lim}}\nolimits}
\def\colim{\mathop{\mathrm{colim}}\nolimits}
\def\Spec{\mathop{\mathrm{Spec}}}
\def\Hom{\mathop{\mathrm{Hom}}\nolimits}
\def\Ext{\mathop{\mathrm{Ext}}\nolimits}
\def\SheafHom{\mathop{\mathcal{H}\!\mathit{om}}\nolimits}
\def\SheafExt{\mathop{\mathcal{E}\!\mathit{xt}}\nolimits}
\def\Sch{\mathit{Sch}}
\def\Mor{\mathop{\mathrm{Mor}}\nolimits}
\def\Ob{\mathop{\mathrm{Ob}}\nolimits}
\def\Sh{\mathop{\mathit{Sh}}\nolimits}
\def\NL{\mathop{N\!L}\nolimits}
\def\CH{\mathop{\mathrm{CH}}\nolimits}
\def\proetale{{pro\text{-}\acute{e}tale}}
\def\etale{{\acute{e}tale}}
\def\QCoh{\mathit{QCoh}}
\def\Ker{\mathop{\mathrm{Ker}}}
\def\Im{\mathop{\mathrm{Im}}}
\def\Coker{\mathop{\mathrm{Coker}}}
\def\Coim{\mathop{\mathrm{Coim}}}

% Boxtimes
%
\DeclareMathSymbol{\boxtimes}{\mathbin}{AMSa}{"02}

%
% Macros for moduli stacks/spaces
%
\def\QCohstack{\mathcal{QC}\!\mathit{oh}}
\def\Cohstack{\mathcal{C}\!\mathit{oh}}
\def\Spacesstack{\mathcal{S}\!\mathit{paces}}
\def\Quotfunctor{\mathrm{Quot}}
\def\Hilbfunctor{\mathrm{Hilb}}
\def\Curvesstack{\mathcal{C}\!\mathit{urves}}
\def\Polarizedstack{\mathcal{P}\!\mathit{olarized}}
\def\Complexesstack{\mathcal{C}\!\mathit{omplexes}}
% \Pic is the operator that assigns to X its picard group, usage \Pic(X)
% \Picardstack_{X/B} denotes the Picard stack of X over B
% \Picardfunctor_{X/B} denotes the Picard functor of X over B
\def\Pic{\mathop{\mathrm{Pic}}\nolimits}
\def\Picardstack{\mathcal{P}\!\mathit{ic}}
\def\Picardfunctor{\mathrm{Pic}}
\def\Deformationcategory{\mathcal{D}\!\mathit{ef}}
\title{Factorization of the zero scheme of a section of a tensor product of invertible sheaves over a flat family of relative dimension one}
\author{The Stacks Project Authors}
\date{stacks-project:chow.tex:local-snapshot-20261012; GNU FDL 1.2 or later}
\begin{document}
\maketitle
\noindent\textit{Editorial note.} All delivered mathematical text is the author's own, copied from the pinned Stacks Project chapter file chow.tex; no statement or proof step was rewritten, reordered or removed.\par\medskip
\noindent\textit{Editorial note.} Preamble provenance: the pinned chapter file opens with \textbackslash{}input\{preamble\} and carries no preamble text of its own, so the project's single Stacks-wide preamble staged at sources/stacks-chow/preamble.tex is excerpted verbatim as two spans, lines 14-18 (Xy-pic setup) and lines 154-217 (the native theorem declarations and mathematical macros). Nothing was added, removed or redefined; the wrapper class and the amsmath/amssymb/amsthm/graphicx base packages are supplied by the assembler.\par\medskip
\noindent\textit{Editorial note.} Same-chapter pointers are located at their \textbackslash{}label lines inside the pinned chapter file; cross-chapter pointers are located in the project's own tag list staged in the snapshot (sources/stacks-chow/tags.txt), at the verbatim lines carrying their canonical tags.\par\medskip
\noindent\textit{Editorial note.} The counted proof opens with a displayed commutative diagram in Xy-pic, which is why the native xy package of the shared Stacks preamble is retained; the diagram is the author's own and was not redrawn.\par\medskip
\noindent\textit{Allowed standard prerequisites (editorial).}\begin{itemize}
\item Relative spectrum of a quasi-coherent sheaf of algebras. Same chapter, chow.tex, Section 'Relative spectrum' (Stacks tag 01LQ, label section-relative-spectrum). Identifies the relative spectrum of the graded quotient algebra A as the flat family p : X -> Y carrying the two sections s\_i; the actual construction of A and of the adjoint sections is inside the counted span.
\item Zero scheme of a global section of an invertible module. Divisors, tag 01WX (label lemma-section-and-zero-scheme) and Schemes, tag 01VZ on closed subschemes cut out by an ideal sheaf. Supplies that D\_i, D and the closed subschemes cut out by the local equations z\_i are the zero schemes of the corresponding sections; the identification of D\_i with the Spec of Sym of the dual invertible module is carried out in the counted span.
\end{itemize}\medskip
\begin{lemma}
\label{lemma-decompose-section}
Let $Y$ be a scheme. Let $\mathcal{L}_i$, $i = 1, 2$ be invertible
$\mathcal{O}_Y$-modules. Let $s$ be a global section of
$\mathcal{L}_1 \otimes_{\mathcal{O}_Y} \mathcal{L}_2$.
Denote $i : D \to Y$ the zero scheme of $s$.
Then there exists a commutative diagram
$$
\xymatrix{
D_1 \ar[r]_{i_1} \ar[d]_{p_1} &
L \ar[d]^p &
D_2 \ar[l]^{i_2} \ar[d]^{p_2} \\
D \ar[r]^i &
Y &
D \ar[l]_i
}
$$
and sections $s_i$ of $p^*\mathcal{L}_i$ such that
the following hold:
\begin{enumerate}
\item $p^*s = s_1 \otimes s_2$,
\item $p$ is of finite type and flat of relative dimension $1$,
\item $D_i$ is the zero scheme of $s_i$,
\item $D_i \cong
\underline{\Spec}(\text{Sym}^*(\mathcal{L}_{3 - i}^{\otimes -1})|_D))$
over $D$ for $i = 1, 2$,
\item $p^{-1}D = D_1 \cup D_2$ (scheme theoretic union),
\item $D_1 \cap D_2$ (scheme theoretic intersection) maps
isomorphically to $D$, and
\item $D_1 \cap D_2 \to D_i$
is the zero section of the line bundle $D_i \to D$ for $i = 1, 2$.
\end{enumerate}
Moreover, the formation of this diagram and the sections $s_i$
commutes with arbitrary base change.
\end{lemma}

\begin{proof}
Let $p : X \to Y$ be the relative spectrum of the quasi-coherent
sheaf of $\mathcal{O}_Y$-algebras
$$
\mathcal{A} =
\left(\bigoplus\nolimits_{a_1, a_2 \geq 0}
\mathcal{L}_1^{\otimes -a_1} \otimes_{\mathcal{O}_Y}
\mathcal{L}_2^{\otimes -a_2}\right)/\mathcal{J}
$$
where $\mathcal{J}$ is the ideal generated by local sections of
the form $st - t$ for $t$ a local section of any summand
$\mathcal{L}_1^{\otimes -a_1} \otimes \mathcal{L}_2^{\otimes -a_2}$
with $a_1, a_2 > 0$. The sections $s_i$ viewed as maps
$p^*\mathcal{L}_i^{\otimes -1} \to \mathcal{O}_X$ are defined as the adjoints
of the maps $\mathcal{L}_i^{\otimes -1} \to \mathcal{A} = p_*\mathcal{O}_X$.
For any $y \in Y$ we can choose an affine
open $V \subset Y$, say $V = \Spec(A)$, containing $y$ and
trivializations $z_i : \mathcal{O}_V \to \mathcal{L}_i^{\otimes -1}|_V$.
Observe that $f = s(z_1z_2) \in A$ cuts out the closed subscheme $D$.
Then clearly
$$
p^{-1}(V) = \Spec(B[z_1, z_2]/(z_1 z_2 - f))
$$
Since $D_i$ is cut out by $z_i$ everything is clear.
\end{proof}

\end{document}
